\documentclass[10pt,a4paper]{amsart}
\usepackage[margin=19mm]{geometry}
\usepackage{amsmath,amssymb,mathtools}
\usepackage[scr=rsfs,cal=euler]{mathalfa}
\usepackage{xfrac}
\usepackage[unicode,hidelinks,bookmarks=false]{hyperref}
\newcommand{\ie}{i.e.\ }
\newcommand{\cf}{cf.\ }
\newcommand{\resp}{respectively\ }
\newcommand{\Subsection}{\subsection}
\providecommand{\nobreakdash}{\nobreak\hbox{-}}
\setlength{\emergencystretch}{2em}
\multlinegap0pt
\usepackage{bm}
\usepackage{amssymb}
\usepackage[shortlabels]{enumitem}
\usepackage{mathtools}
\newtheorem{theorem}{Theorem}
\title{\Large Neural Network Operator-Based Fractal Approximation: Smoothness Preservation and Convergence Analysis}
\author{Aaqib Ayoub Bhat, Asif Khan,, M. Mursaleen}
\date{Source: arXiv:2505.06229v2 (CC BY 4.0)}
\begin{document}
\maketitle

\noindent\emph{Source and licence.} \Large Neural Network Operator-Based Fractal Approximation: Smoothness Preservation and Convergence Analysis. Version arXiv:2505.06229v2 (CC BY 4.0), distributed by arXiv under the Creative Commons Attribution 4.0 International licence, http://creativecommons.org/licenses/by/4.0/, as recorded in the arXiv licence metadata for this exact version and shown on the abstract page https://arxiv.org/abs/2505.06229v2. Full licence notice: \texttt{licenses/arxiv-2505.06229-CC-BY-4.0.txt}.\par\smallskip

\noindent\textit{Editorial scope.} Counted unit: the article's theorem Thml (source0.tex lines 534--546). The article's complete original proof of it is delivered with it (source0.tex lines 548--582, one \texttt{proof} environment of 35 lines). No mathematics is written, repaired or re-derived: the statement and the proof are the article's own lines, in the article's own notation, and no retained line is substituted or re-flowed.  Retained uncounted as the source-local context the proof needs: lines 229--231; lines 233--235.  Nothing was omitted from any retained span, so the package carries no omission record.  No retained line contains a citation, so the wrapper carries no bibliography and no bibliography entry.  No retained line contains a figure environment, an image-inclusion command or drawing code, so no image is delivered and the package declares no asset dependency.  Editorial formatting only, in addition: an AMS \texttt{amsart} preamble that restates the article's own declarations and macros used by the retained lines, namely the environments \texttt{theorem} on separate counters, together with the standard packages those lines need.  The retained environments print in this document as Theorem 1 in order of appearance, whereas the article prints the same items with the numbering of its own full document.  The complete native source of the article as distributed in the arXiv e-print for this exact version is bundled unchanged as \texttt{sources/177877/source0.tex}.
	\begin{equation}\label{defli}
		\pounds_i(x)= \frac{x_i -x_{i-1}}{x_N-x_0}x +\frac{x_N x_{i-1} -x_0 x_i}{x_N-x_0}, \quad \forall i=1,2,\cdots,N
	\end{equation}

	\begin{equation}\label{deffi}
		\mathcal{J}_i(x,y)=\alpha_{i}.y+f(\pounds(x))-\alpha_{i}S_{n,\rho}(f, x), \quad \forall i=1,2,\cdots,N.
	\end{equation}\\

	\begin{theorem}\label{Thml}
		Let $f \in \mathcal{H}^{0,\mu}$. Suppose $\Delta=\left\{a=x_{0}, x_{1}, x_{2}, \ldots, x_{N}=b\right\}$ is a partition of closed interval $[a,b]$ with $x_{0}<x_{1}<\cdots<x_{N}$. For all $i =1,2,\cdots,N$, let $\pounds_{i}$ be given by (\ref{defli}) and $\mathcal{J}_i$ be defined by (\ref{deffi}) with $\rho \in \mathcal{H}^{0,\mu}$ and scaling function
		$\alpha_{i} \in \mathcal{H}^{0,\mu}$. Let the RB-operator $T:\left(\mathcal{H}^{0,\mu},\|\cdot\|_{0, \mu}\right) \rightarrow\left(\mathcal{H}^{0,\mu},\|\cdot\|_{0, \mu}\right)$  be defined by:
		
		\begin{equation*}
			T\phi(x)=f(x)+\boldsymbol{\alpha}_{i}\left(\pounds_{i}^{-1}(x)\right)\left(\phi-S_{n,\rho} f\right)\left(\pounds_{i}^{-1}(x)\right).
		\end{equation*}
		
		\noindent If the scaling vector $\alpha$ satisfies 
		\begin{equation}\label{scale}
			\max \left\{\frac{\left\|\alpha_{i}\right\|_{\infty}}{a_{i}^{\mu}}: i =1,2,\cdots,N\right\}<1,
		\end{equation} then  $T$ admits a unique fixed point $f^{\alpha} \in \mathcal{H}^{0,\mu}$.\\
	\end{theorem}

	\begin{proof}
		
		Consider an arbitrary element $\phi \in \mathcal{H}^{0,\mu}$. Our aim is to show that $T\phi \in \mathcal{H}^{0,\mu} $. For each $n \in \mathbb{N}$,
		\begin{align*}
			|T \phi|_{\mu}= & \sup _{\underset{x \neq y}{x, y \in [a,b],}} \frac{|T \phi(x)-T \phi(y)|}{|x-y|^{\mu}} \\
			= & \sup _{\underset{x \neq y}{x, y \in [x_{i-1},x_i], }} \frac{\left|f(x)-f(y)+\alpha_{i}\left(\phi-S_{n,\rho}(f)\right) \circ\left(\pounds_{i}^{-1}(x)\right)-\alpha_{i}\left(\phi-S_{n,\rho}(f)\right) \circ\left(\pounds_{i}^{-1}(y)\right)\right|}{|x-y|^{\mu}} \\
			\leq & \sup _{\underset{x \neq y}{x, y \in [x_{i-1},x_i], }}\left\{\frac{|f(x)-f(y)|}{|x-y|^{\mu}}\right\}+\max _{i =1,2,\cdots,N}\left(\left\|\alpha_{i}\right\|_{\infty}\right) \sup _{\underset{x \neq y}{x, y \in [x_{i-1},x_i], }}\left[\frac{\left|\phi\left(\pounds_{i}^{-1}(x)\right)-\phi\left(\pounds_{i}^{-1}(y)\right)\right|}{|x-y|^{\mu}}\right. \\
			& \left.+\frac{\left|S_{n,\rho}\left(f ; \pounds_{i}^{-1}(x)\right)-S_{n,\rho}\left(f ; \pounds_{i}^{-1}(y)\right)\right|}{|x-y|^{\mu}}\right]
		\end{align*}
		
		
		From assumptions $f \in \mathcal{H}^{0,\mu},~ S_{n,\rho}(f) \in \mathcal{H}^{0,\mu}$, the above inequality implies $|T \phi|_{\mu}<\infty$, and hence $T \phi \in \mathcal{H}^{0,\mu}$.\\
		Our next claim is to prove that $T$ is contraction. For this, let $\phi,~ \psi \in \mathcal{H}^{0,\mu}$, we have
		
		\begin{equation}\label{10}
			\left\|T\phi - T\psi\right\|_{\infty} \leq \max _{i =1,2,\cdots,N}\left(\left\|\alpha_{i}\right\|_{\infty}\right)\left\|T\phi - T\psi\right\|_{\infty}
		\end{equation}
		
		
		Following a similar approach as in the estimation of $|T g|_{\mu}$, we derive
		
		\begin{equation}\label{11}
			\left|T\phi - T\psi\right|_{\mu} \leq \max _{i =1,2,\cdots,N}\left(\frac{\left\|\alpha_{i}\right\|_{\infty}}{a_{i}^{\mu}}\right)\left|\phi - \psi\right|_{\mu}
		\end{equation}
		
		From the inequalities (\ref{10}) and (\ref{11}), we get
		\begin{align}\label{contr}
			\left\|T\phi - T\psi\right\|_{0, \mu} & =\max \left\{\left\|T\phi - T\psi\right\|_{\infty},\left|T\phi - T\psi\right|_{\mu}\right\}\nonumber \\
			& \leq \max \left\{\max _{i =1,2,\cdots,N}\left(\left\|\alpha_{i}\right\|_{\infty}\right)\left\|\phi - \psi\right\|_{\infty}, \max _{i =1,2,\cdots,N}\left(\frac{\left\|\alpha_{i}\right\|_{\infty}}{a_{i}^{\mu}}\right)\left|\phi - \psi\right|_{\mu}\right\} \nonumber \\
			& =\max _{i =1,2,\cdots,N}\left(\frac{\left\|\alpha_{i}\right\|_{\infty}}{a_{i}^{\mu}}\right) \max \left\{\left\|\phi - \psi\right\|_{\infty},\left|\phi - \psi\right|_{\mu}\right\} \nonumber\\
			& =\max _{i =1,2,\cdots,N}\left(\frac{\left\|\alpha_{i}\right\|_{\infty}}{a_{i}^{\mu}}\right)\left\|\phi - \psi\right\|_{0, \mu} .
		\end{align}
		
		Hence, from assumption (\ref{scale}) and from inequality given in (\ref{contr}), the operator $T$ is a contraction defined on a complete metric space $\left(\mathcal{H}^{0,\mu}, \|\cdot\|_{0, \mu}\right)$. Therefore, by Banach contraction principle, $T$ admits a unique fixed point, say $f_{n}^{\alpha}$. \\
	\end{proof}

\end{document}
